\documentclass[11pt]{article}
\usepackage[margin=1.2in, top=1in, bottom=1in]{geometry}
\usepackage{amsmath,amssymb}
\usepackage{microtype}
\usepackage{lmodern}
\usepackage[T1]{fontenc}
\usepackage[colorlinks=true,linkcolor=black]{hyperref}

\pagestyle{empty}

\begin{document}

\begin{center}
  {\large\bfseries RetroDiff: Training Objective Derivation}
\end{center}

\medskip

\noindent
Let $x$ denote a clean data sample, $x_t$ its noised version at diffusion step $t$,
$q$ the fixed forward noising distribution, and $p_\theta$ the learned reverse distribution
parameterised by $\theta$.  Starting from the log-likelihood, Jensen's inequality gives the
\emph{evidence lower bound} (ELBO):

\begin{align}
  \log p_\theta(x) \;\geq\;
  \mathbb{E}_{q(x_{1:T}|x)}\!\left[\log\frac{p_\theta(x_{0:T})}{q(x_{1:T}|x)}\right]
  = -\mathcal{L}_{\mathrm{ELBO}}
  \label{eq:elbo}
\end{align}

\noindent
Maximising the right-hand side of \eqref{eq:elbo} is equivalent to minimising
$\mathcal{L}_{\mathrm{ELBO}}$.  Expanding the joint distributions and applying the Markov
structure of the forward process yields the canonical decomposition into a reconstruction
term and $T\!-\!1$ per-step KL divergences:

\begin{align}
  \mathcal{L}_{\mathrm{ELBO}}
  &= \underbrace{\mathbb{E}_{q}\!\left[-\log p_\theta(x_0\mid x_1)\right]}_{\text{reconstruction}}
   + \sum_{t=2}^{T}
     \underbrace{\mathrm{KL}\!\left(
       q(x_{t-1}\mid x_t,x)
       \;\|\;
       p_\theta(x_{t-1}\mid x_t)
     \right)}_{\text{step-}t\text{ denoising}}
  \label{eq:decomp}
\end{align}

\noindent
Each KL term in \eqref{eq:decomp} has a tractable Gaussian form because the forward
posterior $q(x_{t-1}|x_t,x)$ is Gaussian.  \textbf{RetroDiff}'s key contribution is to
condition the objective on a set $R(x)$ of $k$ retrieved exemplars drawn from the training
corpus at inference time, leading to the retrieval-augmented objective:

\begin{align}
  \mathcal{L}_{\mathrm{RetroDiff}}(\theta)
  = \mathbb{E}_{x\sim p_{\mathrm{data}}}\!\left[
      \mathcal{L}_{\mathrm{ELBO}}\!\left(x \;\big|\; R(x)\right)
    \right]
  \label{eq:retro}
\end{align}

\noindent
Minimising \eqref{eq:retro} trains the reverse process to exploit retrieved neighbours,
enabling sharper reconstructions without architectural changes.

\end{document}
